\documentclass[14pt,fleqn]{extarticle}
\RequirePackage{prepwell-eng}

\newcommand\cmps{\text{ cm/s}}
\previewoff 

\begin{document} 
\begin{snippet}
    
    \correct
    
    A cube with sides 1m is filled with water, and has a tiny hole through 
    which water drains into a cylinder of radius 1m\newline 
    
    If the water level in the cube is falling at the rate of 1 cm/s, then 
    it rising in the cylinder at the rate of $\frac{1}\pi$ cm/s
    
    \reason
    
    The situation is as shown below 
    \begin{center}
\includegraphics[scale=1.2]{116-A.svg}
\end{center}

The rate at which the \underline{cube loses} water is the same as the rate which 
the \underline{cylinder gains} water. And therefore 
\begin{align}
	L^2\cdot\frac{dh}{dt} &= \pi R^2\cdot \frac{dH}{dt} \\
	\left(1\text{m} \right)^2\cdot 1\cmps &= \pi \left(1\text{m} \right)^2\cdot\frac{dH}{dt} \\
	\therefore \frac{dH}{dt} &= \frac{1}\pi\cmps 
\end{align}
    
\end{snippet} 
\end{document} 